%% Part 2.2 calculi fragment: v1.1 Part III, lines 3927--4305, plus
%% D-P3-001, the first-order half of D-P3-002, and D-P3-011.
%% Rules occur only in the registry sheets named below.  Apparatus: A-037.

\chapter{First-Order Calculi}
\label{sec:p02-first-order-calculi}

The active first-order proof systems keep their languages and equality
standings explicit.  Natural deduction uses
\hyperref[sheet:nd-fol]{\texttt{sheet:nd-fol}}.  The equality-free sequent
calculus uses \hyperref[sheet:lk]{\texttt{sheet:lk}}, while
\hyperref[sheet:lk-eq]{\texttt{sheet:lk-eq}} is the distinct equality
extension.  The equality-free Hilbert presentation uses
\hyperref[sheet:hilbert-fol]{\texttt{sheet:hilbert-fol}}.  Analytic tableaux
use \hyperref[sheet:tableaux-fol]{\texttt{sheet:tableaux-fol}}.  No rule
stated on those sheets is silently imported into another calculus.

\section{Natural deduction and sequents}
\label{sec:p02-fol-nd-lk}

\begin{hbdefinition}{def:p02-fol-derivability}{First-order natural-deduction derivability}
For \(\Gamma\subseteq\mathsf{Fm}(\mathcal L)\),
\index{derivability!first-order natural deduction}
\[
 \Gamma\derives_{\mathsf{ND}_{\mathsf{FOL}}}\varphi
\]
means that a finite derivation generated by
\hyperref[sheet:nd-fol]{\texttt{sheet:nd-fol}} concludes \(\varphi\) and
every open assumption belongs to \(\Gamma\).  The discharge labels and
eigenvariables are metalinguistic data.  Every carried discharge label marks
at least one open occurrence; the calculus has neither a repetition rule nor
vacuous discharge.  This definition names derivability in that sheet, not an
unnamed universal proof relation.
\end{hbdefinition}

\begin{hblemma}{thm:p02-reduct-invariance}{proved=annex:p2.fol.calc.1}{Invariance under reduct}
Let \(\mathcal L'\) extend \(\mathcal L\), let \(\mathcal M'\) be an
\(\mathcal L'\)-structure, and let \(s:\mathsf{Var}\to
\dom\mathcal M'\) be an assignment.  For every \(\mathcal L\)-term \(t\)
and \(\mathcal L\)-formula \(\varphi\),
\[
 \operatorname{val}^{\mathcal M'}_s(t)
 =\operatorname{val}^{\mathcal M'\!\upharpoonright\mathcal L}_s(t),
 \qquad
 \mathcal M',s\entails\varphi
 \quad\text{iff}\quad
 \mathcal M'\!\upharpoonright\mathcal L,s\entails\varphi.
\]
\end{hblemma}

The source examples of universal reasoning and existential elimination are
now lookup pointers rather than body proofs.  The former uses a variable not
free in either premise before \(\forall I\), followed by the explicit
alpha-conversion rule; the latter confines its fresh witness to the minor
derivation of \(\exists E\).  The complete side conditions are on
\hyperref[sheet:nd-fol]{\texttt{sheet:nd-fol}}.

An \(\mathsf{LK}\) derivation is a finite tree of finite-multiset sequents
generated by the chosen registry sheet.  Sequent satisfaction keeps the
Part-2.1 reading: the antecedent is conjunctive, the succedent disjunctive,
and empty antecedent/succedent have the values \(\topv\)/\(\bottom\).
The equality-free and equality-bearing sheets are intentionally distinct.

\begin{hbtheorem}{thm:p02-fol-nd-soundness}{proved=annex:p2.fol.calc.2}{Soundness of first-order natural deduction}
For every \(\Gamma\cup\{\varphi\}\subseteq\mathsf{Fm}(\mathcal L)\),
\[
 \Gamma\derives_{\mathsf{ND}_{\mathsf{FOL}}}\varphi
 \quad\Longrightarrow\quad
 \Gamma\entails\varphi.
\]
The semantic consequence relation ranges over \(\mathcal L\)-structures with
nonempty domains and equality interpreted as identity.
\end{hbtheorem}

\begin{hbtheorem}{thm:p02-fol-lk-soundness}{proved=annex:p2.fol.calc.3}{Soundness of the first-order sequent calculi}
Every sequent derivable on \hyperref[sheet:lk]{\texttt{sheet:lk}} is
satisfied by every \(\mathcal L\)-structure and assignment.  The same holds
for \hyperref[sheet:lk-eq]{\texttt{sheet:lk-eq}} when equality is interpreted
as actual identity.
\end{hbtheorem}

\section{Equality metatheory}
\label{sec:p02-fol-equality-metatheory}

\begin{hbtheorem}{thm:p02-equality-laws}{proved=annex:p2.fol.calc.4}{Derived equality laws}
For all \(\mathcal L\)-terms \(r,t,u\), every formula \(\varphi\), and every
variable \(x\), the natural-deduction equality rules on
\hyperref[sheet:nd-fol]{\texttt{sheet:nd-fol}} derive:
\[
 \varnothing\derives_{\mathsf{ND}_{\mathsf{FOL}}}t=t,
 \qquad
 \{r=t\}\derives_{\mathsf{ND}_{\mathsf{FOL}}}t=r,
\]
\[
 \{r=t,t=u\}\derives_{\mathsf{ND}_{\mathsf{FOL}}}r=u,
 \qquad
 \{r=t,\varphi[r/x]\}
   \derives_{\mathsf{ND}_{\mathsf{FOL}}}\varphi[t/x].
\]
Every displayed substitution is the total capture-avoiding operation of the
syntax chapter.  Reflexivity and replacement are the registered primitive
rules; symmetry and transitivity are derived, not silently added axioms.
\end{hbtheorem}

The free-standing term-congruence result used by the quotient term model is
Theorem~\ref{thm:p02-equality-congruence}.  Together with the preceding
derived laws, it supplies the consolidated equality metatheory required by
D-P3-011.  The symmetric replacement instance for
\hyperref[sheet:lk-eq]{\texttt{sheet:lk-eq}} is likewise derived by symmetry
and Cut; the sheet comments record that it is not primitive.

\section{An equality-free Hilbert calculus}
\label{sec:p02-hilbert-fol}

\begin{hbdefinition}{def:p02-hilbert-fol-dependencies}{Hilbert derivation with dependencies}
Let \(\mathsf H_{\mathsf{FOL}}\) be
\hyperref[sheet:hilbert-fol]{\texttt{sheet:hilbert-fol}}, whose language is
the equality-free \(\{\neg,\to,\forall,\exists\}\)-fragment.  A
\emph{dependency-annotated derivation from \(\Gamma\)}
\index{Hilbert calculus!first-order}
\index{derivation!dependency-annotated} is a finite sequence
of pairs \((\chi_i,D_i)\), where \(D_i\subseteq\Gamma\) is finite, with
exactly these clauses:
\begin{enumerate}
  \item if \(\chi_i\in\Gamma\), then \(D_i=\{\chi_i\}\);
  \item an axiom-schema instance has \(D_i=\varnothing\);
  \item modus ponens from lines \(j,k<i\) has
        \(D_i=D_j\cup D_k\); and
  \item generalization from \((\chi_j,D_j)\) to
        \((\forall x\,\chi_j,D_j)\) is permitted exactly when
        \(x\notin\FV(\delta)\) for every \(\delta\in D_j\).
\end{enumerate}
Write \(\Gamma\derives_{\mathsf H_{\mathsf{FOL}}}\chi\) iff some such
sequence ends in \((\chi,D)\).  The annotation is metalinguistic and
eliminable: erasing it gives the sheet's finite formula sequence, while
reconstructing it records exactly which premises license each generalization.
It is needed below because the first-order deduction theorem fails for an
unrestricted generalization rule.
\end{hbdefinition}

\begin{hbtheorem}{thm:p02-hilbert-fol-deduction}{proved=annex:p2.fol.calc.5}{Hilbert-FOL deduction theorem}
For every set \(\Gamma\) and formulas \(\varphi,\psi\) of the sheet's
fragment,
\[
 \Gamma\cup\{\varphi\}
   \derives_{\mathsf H_{\mathsf{FOL}}}\psi
 \quad\Longleftrightarrow\quad
 \Gamma\derives_{\mathsf H_{\mathsf{FOL}}}(\varphi\to\psi).
\]
The derivability relation is the dependency-annotated one just defined; the
theorem does not license generalizing a variable free in a premise on which
the generalized line depends.
\end{hbtheorem}

This calculus, its quantifier schemas, and its dependency-sensitive
generalization convention are standard Hilbert first-order machinery; see
\cite{Mendelson1997,Enderton2001}.  Its equality-free fragment must not be
reported as literally the same language as the equality-bearing natural-
deduction and sequent sheets.

\section{First-order analytic tableaux}
\label{sec:p02-tableaux-fol}

\begin{hbdefinition}{def:p02-fol-tableau}{First-order tableau}
For a finite set \(\Gamma\) of equality-free \(\mathcal L\)-sentences, let
\(\mathcal L_{\Gamma}\) contain exactly the nonlogical symbols occurring in
\(\Gamma\), and obtain \(\mathcal L_{\Gamma}^*\) by adjoining new parameter
constants \(a_n\) for \(n\in\mathbb N\).  The resulting signature has only
countably many closed terms; fix an enumeration of all of them.  A
\emph{first-order tableau rooted at \(\Gamma\)}\index{tableau!first-order}
is a finitely branching
retained-formula tree generated by
\hyperref[sheet:tableaux-fol]{\texttt{sheet:tableaux-fol}}.  Propositional
and witness-producing occurrences expand once.  A universal occurrence (or
the negation of an existential) records each closed
\(\mathcal L_{\Gamma}^*\)-term already used and remains available at every
unused term.  A branch is \emph{fair} iff every such occurrence is eventually
instantiated with every term in the fixed enumeration, whether or not that
term previously appeared on the branch.  A branch is \emph{closed} iff it
contains \(\bottom\) or
both \(\varphi\) and \(\neg\varphi\) for some \(\varphi\), and is \emph{open}
otherwise; it is \emph{saturated} iff no expansion applicable on that branch
is omitted.  These occurrence marks and parameters are
metalinguistic proof-search data, not new symbols of the original language.
\end{hbdefinition}

\begin{hbtheorem}{thm:p02-fol-tableau-adequacy}{cited=Smullyan1968}{First-order tableau adequacy and countermodel extraction}
Let \(\Gamma\) be a finite set of equality-free \(\mathcal L\)-sentences.
Then \(\Gamma\) is unsatisfiable iff a finite closed tableau rooted at
\(\Gamma\) exists.  Moreover, for every fair saturated open branch \(B\),
define its positive atomic part by
\[
 B^{+}_{\mathrm{at}}
 \defeq
 \{R(t_1,\ldots,t_n):R(t_1,\ldots,t_n)\text{ occurs on }B\}.
\]
The corresponding Herbrand structure is
\(\mathcal M_{B^{+}_{\mathrm{at}}}\).  Its domain is the closed
\(\mathcal L_{\Gamma}^*\)-terms, its functions act by term formation, and
its relations are exactly the members of \(B^{+}_{\mathrm{at}}\).
The branch truth induction gives
\[
 \mathcal M_{B^{+}_{\mathrm{at}}}\entails B.
\]
Hence its
\(\mathcal L_{\Gamma}\)-reduct models \(\Gamma\), and interpreting every
unused \(\mathcal L\)-symbol arbitrarily expands it to an
\(\mathcal L\)-model of \(\Gamma\).
This is the countermodel-extraction clause, not merely a closure test.
\end{hbtheorem}

The equality-free scope is load-bearing.  Identity tableaux require an
additional equality calculus; neither the equality rules of
\hyperref[sheet:nd-fol]{\texttt{sheet:nd-fol}} nor those of
\hyperref[sheet:lk-eq]{\texttt{sheet:lk-eq}} are silently available
on the tableau sheet.  Together with the already verified propositional entry,
this completes D-P3-002.

\begin{landscape}
\section{First-order proof-system comparison}
\label{sec:p02-fol-calculi-comparison}

\begingroup\scriptsize\setlength{\tabcolsep}{3pt}
\begin{longtable}{@{}p{.65in}p{1.35in}p{1.95in}p{1.90in}p{1.25in}@{}}
\caption{First-order proof systems side by side (apparatus A-037).}
\label{tab:p02-fol-proof-systems}\\
\toprule
System & Rule families & Exact side conditions & Adequacy evidence & Scope and registry home \\ \midrule
\endfirsthead
\toprule
System & Rule families & Exact side conditions & Adequacy evidence & Scope and registry home \\ \midrule
\endhead
Hilbert & Propositional schemas A1--A3; quantifier schemas
\(\forall1,\forall2,\exists1,\exists2\); MP and Gen. & All substitutions are
total and capture-avoiding; in \(\forall2\), \(x\notin\FV(\varphi)\); in
\(\exists2\), \(x\notin\FV(\psi)\); Gen requires \(x\) not free in any
premise on which the line depends. & Soundness and strong completeness are
external cited facts for this standard presentation \cite{Mendelson1997}; the
deduction theorem is proved here as
Theorem~\ref{thm:p02-hilbert-fol-deduction}. & Equality-free
\(\neg,\to,\forall,\exists\) fragment;
\hyperref[sheet:hilbert-fol]{\texttt{sheet:hilbert-fol}}. \\
Natural deduction & Propositional introduction/elimination rules, RAA,
\(\forall I/E\), \(\exists I/E\), alpha-conversion, and \(=I/E\). &
\(\forall I\): \(x\notin\FV(\Gamma)\); \(\exists E\): the witness variable is
absent from the conclusion, all other open assumptions, and the existential
major premise; every discharge is nonvacuous; substitution is total and
capture-avoiding. & Proved soundness,
Theorem~\ref{thm:p02-fol-nd-soundness}, and proved strong completeness,
Theorems~\ref{thm:p02-fol-completeness} and
\ref{thm:p02-fol-completeness-countable}. & Primitive identity;
\hyperref[sheet:nd-fol]{\texttt{sheet:nd-fol}}. \\
\(\mathsf{LK}\) & Contexted initials; structural and propositional left/right
rules; Cut and Mix; \(\forall L/R\), \(\exists L/R\); the equality extension
adds reflexivity and left replacement. & \(\forall R\) and \(\exists L\) use
an eigenvariable free in no context or principal formula; \(\forall L\) and
\(\exists R\) use total capture-avoiding substitution; Mix has positive
multiplicities. & Proved soundness,
Theorem~\ref{thm:p02-fol-lk-soundness}; honest completeness sketch,
Theorem~\ref{thm:p02-fol-lk-completeness}. & Equality-free
\hyperref[sheet:lk]{\texttt{sheet:lk}}; identity extension
\hyperref[sheet:lk-eq]{\texttt{sheet:lk-eq}}. \\
Tableau & Propositional linear/branching expansions; repeated gamma and
one-use delta expansions; complementary-pair or \(\bottom\) closure. & Finite
equality-free sentence root; gamma fairness ranges over every enumerated
closed term; delta uses the least parameter not earlier on the branch; every
substitution is total and capture-avoiding. & Cited finite closed-refutation
completeness and fair-open-branch model extraction,
Theorem~\ref{thm:p02-fol-tableau-adequacy}. & Equality-free;
\hyperref[sheet:tableaux-fol]{\texttt{sheet:tableaux-fol}}. \\
\bottomrule
\end{longtable}
\endgroup
\end{landscape}

\section{Rule-boundary warnings}
\label{sec:p02-fol-rule-warnings}

\begin{hbwarning}{def:p02-fol-quantifier-errors}{Quantifier-rule errors}
None of the following is licensed by the named sheets: generalizing \(P(x)\)
while it remains an open assumption; replacing \(\exists xP(x)\) by a
permanent new premise \(P(c)\); letting the eigenvariable of \(\exists E\)
occur in the conclusion; reversing
\(\forall x\exists y\,R(x,y)\) to
\(\exists y\forall x\,R(x,y)\); or using textual substitution that permits
capture.  The first and third violate freshness conditions, the second leaks
a local witness, the fourth reverses witness dependence, and the fifth
changes the formula's binding structure.
\end{hbwarning}

Source exercises are not retained as exercises.  Equality transitivity is
promoted to Theorem~\ref{thm:p02-equality-laws}; cut-free sequent derivations
belong to Part~7; and capture, eigenvariable, and translation drills are
represented by the definitions and warning above.
